\chapter{Computation and experimental design}\label{sec:methods}
\section{Independent implementations and validation}\label{sec:validation}
The code separates simulation from experiments, analysis, plotting and report generation. A reference routine scans the grid and performs one elementary toppling at a time. The production routine uses a queue of active sites, performing $\lfloor h(x)/4\rfloor$ consecutive legal topplings at each visit and adding that count to the odometer. A circular queue holds at most one pending entry per site. Timestamps identify newly toppled sites without clearing the whole grid after every event.

Heights, counts and mass arithmetic use signed 64-bit integers with a conservative overflow bound. A scalar loop or \texttt{add.at} accumulates repeated input coordinates. Each event's recorded mass is checked against its input and loss bookkeeping. The kernel independently checks actual grid mass through histograms every 128 input grains, and the wrapper sums the grid at each chunk end. In separate validation runs, an independent array-shift simulator checks actual mass after every event. A safety-limit exception aborts the chunk and preserves the stable checkpoint for recovery.

The automated suite has 22 test methods. It covers all $256\times4=1,024$ stable $2\times2$ additions, 1,000 bounded random configurations in each of the Python and compiled cross-checks, open rectangular and degenerate grids, repeated topplings, all four observables, and the corrected handout example. Forward, reverse and independent random legal schedules agree exactly. Extension tests check cylindrical seams, sinks and distances, the sparse benchmark, full-support drive normalization and fixed-seed source frequencies, batch duplicates, odometer sums and area unions. Exact enumeration finds the 192 recurrent $2\times2$ states and verifies the uniform invariant measure for uniform and unequal positive source probabilities, including the fourth-power transition.

Uninterrupted and restarted runs produce identical records, final grids and PCG64 state for every condition, including a failure injected after the metadata commit and before the checkpoint commit. Checkpoints retain the stage, progress, stable grid and RNG state; chunks are committed atomically and hashed. Tests also check stored data types, inclusive-CCDF plotting, statistical normalization, known discrete distributions, sparse samples and autoregressive dependence. Fits use the best converged multistart solution with an independent projected-gradient check. All 144 aggregate fits pass, and unsuccessful individual starts remain in the records.

The project retains raw chunks, completion metadata, frozen configurations, editable sources, tests and generated figures and tables. \texttt{README.md} lists tested validation, smoke, production/resumption, analysis and build commands. A hashed manifest traces report numbers and graphics to their inputs and analysis settings; the question/rubric matrix maps requirements to this evidence.

\section{Profiling and comparable performance}\label{sec:performance}
Profiling identified repeated neighbour construction; caching removed that cost. A later $L=64$ profile placed most time in the active queue, which was then implemented as a small C kernel called through \texttt{ctypes}. Python handles random inputs and the experiment workflow. Sparse snapshots and a histogram every 128 grains reduce storage. Reach calculations depend on the number of toppled sites and sources; each batch here has four sources.

\begin{table}[htbp]
\centering\small
\begin{tabular}{@{}lrrr@{}}
\toprule
Implementation & Median (ms) & Range (ms) & Speed ratio \\
\midrule
Elementary scan, Python & 552.216 & [531.345, 562.245] & 1.0 \\
Bulk queue, Python & 107.030 & [106.292, 107.353] & 5.2 \\
C queue, one event per call & 95.477 & [94.000, 107.416] & 5.8 \\
C queue, full input sequence & 0.167 & [0.166, 0.205] & 3298.7 \\
\bottomrule
\end{tabular}

\caption{Stabilization timings for identical $L=16$ workloads: \R{benchmarkEvents} presampled inputs from an all-three grid, with \R{benchmarkTopples} elementary topplings. After one full-workload warm-up per engine, each of five timing samples repeats complete workloads for at least 0.2\,s, resetting the grid each time. Times are divided by the number of workloads; all implementations produce identical $S,A,D,\ell$ and final grids. Compilation (\R{compileSeconds}\,s) and first load/warm-up (\R{warmupSeconds}\,s) are measured separately. The range runs from minimum to maximum, and speed ratios use the Python scan median. Timings exclude I/O and input generation.}
\label{tab:performance}
\end{table}

The Python queue is \R{queueSpeed} times faster here (Table~\ref{tab:performance}), combining active scheduling, bulk legal toppling and cached geometry. The C stream processes complete streams and computes observables during site visits, reducing Python validation and call overhead. Per-event C calls expose wrapper overhead. Timing samples accumulate at least 0.2 seconds, with C stream times below one millisecond per workload. Table~\ref{tab:performance} and production timings separate compilation and I/O from these stabilization ratios. The environment used Python 3.12.2, NumPy 1.26.4, SciPy 1.13.0, Matplotlib 3.8.4 and Apple Clang 17 (arm64 macOS, eight logical CPUs).

\section{Pilot, production, and clocks}\label{sec:design}
The independent pilot used $L=64$ for open uniform, open biased and cylindrical input, each starting from empty, random stable and all-three grids with separate streams. It followed 60$L^2$ grains in the uniform and cylindrical cases and 100$L^2$ under bias. A separate equilibrated, matched batch pilot checked 81,920 future grains. We then froze the production settings: five seeds per condition, baseline sizes $L=8,16,32,64$, extension sizes $L=32,64$, and 200,000 measured grains per seed. Burn-in is $20L^2$ for baseline and cylinder, and $40L^2$ for bias. Batches inherit the matched baseline's grid after burn-in.

The 50 primary runs contain \R{productionGrains} measured grains and \R{productionEvents} events, plus 2,080,000 burn grains. Each batch run has 50,000 events and copies its stable start and future inputs from the baseline control. Those two million batch grains form paired comparisons. Pilot, production, scheduling, the 100 independent $L=32$ ensemble chains and block resampling use separate random streams. Production took \R{productionWall}\,s wall time, including \R{kernelSeconds}\,s in the kernel and wrapper and \R{ioSeconds}\,s in timed chunk I/O. These measurements include burn-in; the separate diagnostics and analysis have their own costs.

\section{Uncertainty and distribution fitting}\label{sec:statistics}
We estimate primary 95\% Student-$t$ intervals from five independent seed estimates, with four degrees of freedom. Each seed is one replicate, which accounts for dependence within a run but gives limited precision with five replicates. Independent extension differences use Welch intervals; batch comparisons use exact pairing. The time-series comparison uses 24 complete contiguous blocks of 8,192 grains, covering 196,608 grains. The remaining grains are retained in all-event summaries. Block lengths of 4,096 and 16,384 grains and 200 block-bootstrap resamples assess sensitivity. Autocorrelation is calculated separately for height, topplings, loss, dissipation indicators and height frequencies. The initial-positive-pair time estimate uses a finite lag range and has a floor of one, which short-lag anticorrelation can reach.

Fits use the exact counts of raw integer observations, with normalized likelihood
\begin{equation}
 p(x\mid\alpha,\lambda)=\frac{x^{-\alpha}e^{-\lambda x}}
 {\sum_{n=x_{\min}}^{x_{\max}}n^{-\alpha}e^{-\lambda n}}.
 \label{eq:fit}
\end{equation}
The candidates are power ($\lambda=0$), cutoff power and exponential ($\alpha=0$), fitted to identical observations and support. Numerical bounds are $-10\leq\alpha\leq10$ and $0\leq\lambda x_{\max}\leq10^5$; no fitted free parameter reaches them. Finite support allows negative effective slopes. Each seed must retain at least 1,000 observations. The fixed lower limits are $(16,8,1,3)$ for $(S,A,D,\ell)$; sensitivity checks halve and double them, with a minimum of one. For $S$, the upper fitting limit $L^2$ selects an intermediate range. The other upper limits, $L^2$, $3L^2+1$ and $2L-1$, are physical bounds. We record observed maxima separately from these fitting limits.

Parameter intervals are calculated across the five seed fits. KS distances describe discrepancies between fitted and empirical CDFs. Marginal log-score differences and their block sensitivity use the fitting data, so these comparisons are descriptive; the cutoff model also has an extra parameter. For logarithmic displays, counts are divided by bin width and sample count. Fits use the original integer observations, with each size and model analysed separately.
